\section{Introducción}
En este trabajo, se aborda la modelización de la trayectoria de la nave Orión en la misión Artemis II, mediante la aplicación de métodos numéricos. El objetivo principal es comparar el comportamiento observado del modelo teórico, con el de los experimentos. 

La simulación se realizó con tres métodos de distinto orden, con fin de compararlos y observar el impacto que se genera en los resultados al realizar cambios en los valores iniciales.

El trabajo consta de dos partes. La primera, se centra en la obtención de la órbita de la luna sobre la tierra, y la segunda, en simular la trayectoria de orión en su recorrido entre la tierra y la luna.